% Require in LaTeX preamble: \usepackage{pgfplots}, \usepackage{tikz}
\begin{figure*}[htbp]
  \centering
  \begin{tikzpicture}
    \pgfplotsset{every axis label/.append style={font=\large}, every tick label/.append style={font=\large}, title style={font=\Large\bfseries}}
    % Panel 1: Duration per Step
    \begin{axis}[
      name=ax1,
      width=0.95\textwidth, height=4.8cm,
      title={Step-by-Step Deep Dive --- Test: \texttt{strat\_pc\_F\_win} (Execution Duration)},
      ylabel={Duration (\mu s)},
      grid=major, grid style={dashed,gray!30},
      legend pos=north west, legend style={font=\small}
    ]
      \addplot[very thick, color=gray!80!black, mark=*] coordinates {(1, 2252) (2, 3025) (3, 2536) (4, 2641) (5, 2632) (6, 2924) (7, 2712)};
      \addlegendentry{SEQ}
      \addplot[very thick, color=orange!80!black, mark=*] coordinates {(1, 1092) (2, 442) (3, 440) (4, 403) (5, 365) (6, 427) (7, 387)};
      \addlegendentry{TBB}
      \addplot[very thick, color=blue!70!black, mark=*] coordinates {(1, 6087) (2, 392) (3, 387) (4, 388) (5, 377) (6, 381) (7, 379)};
      \addlegendentry{OMP}
    \end{axis}
    % Panel 2: Active Threads per Step
    \begin{axis}[
      name=ax2, at={(ax1.below min energy point)}, anchor=above north point, yshift=-1.4cm,
      width=0.95\textwidth, height=4.5cm,
      title={2. Active Threads per Step},
      ylabel={Active Threads},
      grid=major, grid style={dashed,gray!30},
      legend pos=north west, legend style={font=\small}
    ]
      \addplot[very thick, color=gray!80!black, mark=square*] coordinates {(1, 1) (2, 1) (3, 1) (4, 1) (5, 1) (6, 1) (7, 1)};
      \addlegendentry{SEQ}
      \addplot[very thick, color=orange!80!black, mark=square*] coordinates {(1, 12) (2, 20) (3, 19) (4, 20) (5, 20) (6, 20) (7, 20)};
      \addlegendentry{TBB}
      \addplot[very thick, color=blue!70!black, mark=square*] coordinates {(1, 17) (2, 19) (3, 18) (4, 20) (5, 20) (6, 20) (7, 20)};
      \addlegendentry{OMP}
    \end{axis}
    % Panel 3: Workload Composition Stacked Area
    \begin{axis}[
      name=ax3, at={(ax2.below min energy point)}, anchor=above north point, yshift=-1.4cm,
      width=0.95\textwidth, height=4.8cm,
      title={3. Workload Composition per Step},
      ylabel={Number of Regions},
      grid=major, grid style={dashed,gray!30},
      legend pos=north west, legend style={font=\small}
    ]
      \addplot[stack plots=y, area legend, fill=blue!40, draw=blue!70!black] coordinates {(1.0, 440.0) (2.0, 457.0) (3.0, 468.0) (4.0, 475.0) (5.0, 480.0) (6.0, 488.0) (7.0, 491.0)} \closedcycle;
      \addlegendentry{Initial Regions (ToProcess)}
      \addplot[stack plots=y, area legend, fill=green!40, draw=green!70!black] coordinates {(1.0, 660.0) (2.0, 794.0) (3.0, 944.0) (4.0, 998.0) (5.0, 1008.0) (6.0, 1044.0) (7.0, 1075.0)} \closedcycle;
      \addlegendentry{Discrete Predecessors}
      \addplot[stack plots=y, area legend, fill=red!40, draw=red!70!black] coordinates {(1.0, 3271.0) (2.0, 4525.0) (3.0, 5319.0) (4.0, 5559.0) (5.0, 5583.0) (6.0, 5653.0) (7.0, 5684.0)} \closedcycle;
      \addlegendentry{Delay Predecessors}
    \end{axis}
    % Panel 4: Memory Sizes
    \begin{axis}[
      name=ax4, at={(ax3.below min energy point)}, anchor=above north point, yshift=-1.4cm,
      width=0.95\textwidth, height=4.5cm,
      title={4. Memory Footprint per Step},
      xlabel={Batch Step (piFilter Call)}, ylabel={Bytes},
      grid=major, grid style={dashed,gray!30},
      legend pos=north west, legend style={font=\small}
    ]
      \addplot[very thick, color=purple!80!black, mark=triangle*] coordinates {(1.0, 120208.0) (2.0, 124760.0) (3.0, 127704.0) (4.0, 129576.0) (5.0, 130912.0) (6.0, 133104.0) (7.0, 133896.0)};
      \addlegendentry{ToProcess Bytes}
      \addplot[very thick, color=teal!80!black, mark=diamond*] coordinates {(1.0, 130984.0) (2.0, 135808.0) (3.0, 142816.0) (4.0, 144800.0) (5.0, 146216.0) (6.0, 148536.0) (7.0, 149376.0)};
      \addlegendentry{SetG Bytes}
    \end{axis}
  \end{tikzpicture}
  \caption{Step-by-step metrics timeline deep dive for test \texttt{strat\_pc\_F\_win}.}
  \label{fig:timeline_strat_pc_f_win}
\end{figure*}